\chapter{正式上路：顺序结构程序设计}

\begin{introduction}
  \item 从第一个程序开始，逐行读懂 C 程序的骨架
  \item 变量：名字、类型、值、地址四件事
  \item 基本数据类型与字面量的写法
  \item \texttt{printf} 与 \texttt{scanf}：程序如何说话、如何倾听
  \item 表达式、语句、注释与 C89 的声明规则
  \item 顺序结构的综合例题与常见错误
\end{introduction}

\chaptermotto{「所有伟大的程序，都是从一行 \texttt{printf} 开始的。」}

上一章我们把舞台搭好了：CPU 执行指令，数据存在内存里，变量有地址，整数用补码、小数用浮点。现在，是时候让程序真正跑起来了。

这一章只讲一件事：\textbf{顺序结构}——程序从上到下，一条语句接一条语句地执行。没有判断，没有循环，就像照着菜谱一步步做菜。听起来简单，但它是后面所有结构的基础。请把这一章的每个例子都亲手敲一遍。

\section{第一个程序：Hello, world!}

\subsection{完整代码}

\begin{codebox}{hello.c —— 第一个 C 程序}
#include <stdio.h>

int main(void)
{
    printf("Hello, world!\n");
    return 0;
}
\end{codebox}

\begin{outbox}{运行结果}
Hello, world!
\end{outbox}

短短 5 行代码，包含了 C 程序的所有基本骨架。我们一行一行拆开看。

\subsection{逐行解剖}

\subsubsection*{第 1 行：\texttt{\#include <stdio.h>}}

这是一条\textbf{预处理指令}。它以 \code{#} 开头，在真正的编译开始之前，由预处理器处理。

\begin{itemize}
  \item \code{stdio} 是 \en{standard input/output}（标准输入输出）的缩写；
  \item \code{.h} 表示头文件（header）；
  \item 这一行的意思是：把 \file{stdio.h} 这个文件的内容“复制粘贴”到当前位置。
\end{itemize}

为什么要包含它？因为 \code{printf} 这个函数的“说明书”就写在 \file{stdio.h} 里。没有这条指令，编译器就不认识 \code{printf}，会报“隐式声明”的警告或错误。

\begin{tipbox}[头文件不是库]
\code{#include <stdio.h>} 只是把\textbf{声明}拿进来，告诉你“有这么个函数，参数长这样”。函数真正的实现（机器指令）在链接阶段才从 C 标准库里接进来。这就是第 1 章讲的“编译”和“链接”的分工。
\end{tipbox}

\subsubsection*{第 3 行：\texttt{int main(void)}}

这是程序的\textbf{入口}。操作系统加载你的程序后，会先找到 \code{main} 函数，从它的第一行开始执行。

\begin{itemize}
  \item \code{int}：这个函数执行完后，会返回一个整数给操作系统；
  \item \code{main}：函数名，固定写法，不能改；
  \item \code{void}：表示这个函数不接收参数。也可以写成 \code{int main()}，但 \code{void} 更明确；
  \item 括号和后面的花括号：函数体，也就是这个函数要执行的语句。
\end{itemize}

\begin{pitfall}[\texttt{void main()} 是错的]
很多老教材写 \code{void main()}，这在标准 C 里是不规范的。请始终写 \code{int main(void)}，并在结尾写 \code{return 0;}。\code{return 0;} 告诉操作系统“程序正常结束”。
\end{pitfall}

\subsubsection*{第 5 行：\texttt{printf("Hello, world!\textbackslash n");}}

这是一条\textbf{函数调用语句}，拆成四部分：

\begin{center}
\begin{tikzpicture}[x=1cm,y=1cm,>=Stealth]
  \node[font=\ttfamily\small] (p) at (0,0) {printf};
  \node[font=\ttfamily\small] (lp) at (1.5,0) {(};
  \node[font=\ttfamily\small,draw=structurecolor!60,fill=structurecolor!6,rounded corners=1mm,inner sep=2pt] (s) at (5.2,0) {"Hello, world!\textbackslash n"};
  \node[font=\ttfamily\small] (rp) at (8.7,0) {);};
  \node[lbl,above=3mm of p] {函数名};
  \node[lbl,above=3mm of s] {格式串（字符串字面量）};
  \node[lbl,above=3mm of rp] {语句结束};
\end{tikzpicture}
\end{center}

\begin{itemize}
  \item \code{printf} 是标准库函数，负责“格式化输出”；
  \item 双引号括起来的部分叫\textbf{字符串字面量}；
  \item \code{\n} 是转义字符，表示“换行”，不是两个普通字符；
  \item 最后的分号 \code{;} 表示一条语句结束。\textbf{C 语言里，语句必须有分号}。
\end{itemize}

\subsubsection*{第 6 行：\texttt{return 0;}}

把整数 0 返回给操作系统。0 通常表示“成功”，非 0 表示“出错”。在 \code{main} 里，\code{return 0;} 是程序正常退出的标志。

\subsection{程序的结构}

把上面的代码抽象一下，一个最简单的 C 程序长这样：

\begin{center}
\begin{tikzpicture}[x=1cm,y=1cm,>=Stealth,node distance=4mm]
  \node[box,minimum width=9cm,align=left,font=\ttfamily\small] (inc) {\#include <stdio.h>};
  \node[box,minimum width=9cm,align=left,font=\ttfamily\small,below=of inc] (main) {int main(void)};
  \node[box,minimum width=9cm,align=left,font=\ttfamily\small,below=of main] (body) {\{ \quad 语句 1；语句 2；…… \quad \}};
  \node[lbl,right=3mm of inc] {预处理指令};
  \node[lbl,right=3mm of main] {函数头};
  \node[lbl,right=3mm of body] {函数体};
\end{tikzpicture}
\end{center}

\begin{c89box}
从第一个程序开始就养成习惯：
\begin{enumerate}
  \item \code{main} 写成 \code{int main(void)}；
  \item 左花括号 \texttt{\{} 单独占一行（也可以跟在函数头后面，但本书统一单独一行）；
  \item 花括号内的语句统一缩进 4 个空格；
  \item 每条语句以分号结尾；
  \item 最后一个 \code{return 0;} 不要忘。
\end{enumerate}
\end{c89box}

\section{变量：给内存里的格子起个名字}

\subsection{变量是什么}

第 1 章说过：变量是内存中一块有名字、有类型、有地址的存储空间。用一张图表示：

\begin{figure}[htbp]
  \centering
  \begin{tikzpicture}[x=1cm,y=1cm]
    \node[draw=structurecolor!70!black,fill=structurecolor!7,rounded corners=2mm,
          minimum width=8.2cm,minimum height=3.4cm] (box) at (0,0) {};
    \node[anchor=west,font=\small] at (-3.8,1.1) {\textbf{名字}：\texttt{a}};
    \node[anchor=west,font=\small] at (-3.8,0.35) {\textbf{类型}：\texttt{int}（4 字节）};
    \node[anchor=west,font=\small] at (-3.8,-0.4) {\textbf{值}：\texttt{10}};
    \node[anchor=west,font=\small] at (-3.8,-1.15) {\textbf{地址}：\texttt{0x1000}};
    \node[font=\Large\ttfamily,color=structurecolor!80!black] at (2.9,0.35) {10};
    \node[font=\scriptsize,color=black!55] at (2.9,-0.45) {占用 4 个字节};
  \end{tikzpicture}
  \caption{一个 \texttt{int} 变量的四个要素：名字、类型、值、地址}
  \label{fig:variable}
\end{figure}

这四件事里，\textbf{名字}和\textbf{类型}是给程序员和编译器看的；\textbf{值}是运行时真正存的东西；\textbf{地址}是内存中的位置。编译之后，变量名通常就不存在了，机器只认地址。

\begin{tipbox}[把变量想成贴了标签的储物格]
储物格有编号（地址）、有大小（类型）、里面放着东西（值），外面贴着标签（名字）。你写 \code{a = 10;}，就是“把 10 放进标签为 \code{a} 的格子里”。
\end{tipbox}

\subsection{声明、赋值、初始化}

\begin{codebox}{declare.c —— 声明与赋值}
#include <stdio.h>

int main(void)
{
    int a;          /* 声明：告诉编译器“有这么个 int 变量” */
    int b = 20;     /* 初始化：声明的同时给初值 */
    int sum;        /* 只声明，暂时没有有意义的值 */

    a = 10;         /* 赋值：把 10 写进 a 的内存 */
    sum = a + b;    /* 用 a 和 b 的值算出结果，写进 sum */

    printf("a = %d, b = %d, sum = %d\n", a, b, sum);
    return 0;
}
\end{codebox}

\begin{outbox}{运行结果}
a = 10, b = 20, sum = 30
\end{outbox}

\begin{itemize}
  \item \textbf{声明}（declaration）：\code{int a;} 告诉编译器“我要一个 \code{int} 类型的变量，名字叫 \code{a}”。
  \item \textbf{赋值}（assignment）：\code{a = 10;} 把右边表达式的值写进左边变量的内存。
  \item \textbf{初始化}（initialization）：\code{int b = 20;} 在声明的同时给初值。
\end{itemize}

\begin{pitfall}[未初始化的局部变量是“垃圾值”]
\code{int sum;} 之后如果没有赋值就直接打印它，输出的是一个不可预测的值。局部变量不会自动清零！这是初学者最常见的 bug 之一。\textbf{变量在使用前一定要先赋值}。
\end{pitfall}

\subsection{C89 的声明规则}

这是本书反复强调的重点。在 C89 里，\textbf{一个语句块内的所有声明必须放在块的开头}，在任何执行语句之前。

\begin{codebox}{order.c —— 正确的声明顺序（C89）}
#include <stdio.h>

int main(void)
{
    int a = 3;
    int b = 4;
    int sum;

    sum = a + b;
    printf("%d\n", sum);
    return 0;
}
\end{codebox}

\begin{codebox}{wrong.c —— 错误示范（C99 才允许）}
#include <stdio.h>

int main(void)
{
    int a = 3;

    printf("a = %d\n", a);

    int b = 4;      /* 错误：声明出现在了执行语句之后 */
    printf("b = %d\n", b);
    return 0;
}
\end{codebox}

在 C89 编译器下，\file{wrong.c} 会报类似这样的错误：

\begin{outbox}{编译错误（GCC，-std=c89）}
wrong.c: In function 'main':
wrong.c:9:5: error: ISO C90 forbids mixed declarations and code
\end{outbox}

\begin{c89box}[正确的写法]
把 \code{int b = 4;} 挪到 \code{printf} 之前、所有执行语句之前：
\begin{codebox}{right.c —— 声明全部提前}
#include <stdio.h>

int main(void)
{
    int a = 3;
    int b = 4;

    printf("a = %d\n", a);
    printf("b = %d\n", b);
    return 0;
}
\end{codebox}
\end{c89box}

\subsection{命名规则与风格}

变量名不是随便起的，必须遵守语法：

\begin{itemize}
  \item 只能包含字母、数字、下划线，且不能以数字开头；
  \item 区分大小写：\code{sum} 和 \code{Sum} 是两个不同的变量；
  \item 不能使用关键字（如 \code{int}、\code{return}、\code{for}）；
  \item 不能和标准库函数重名（例如不要定义叫 \code{printf} 的变量）。
\end{itemize}

风格上，本书推荐：

\begin{center}
\begin{tabular}{@{}lll@{}}
\toprule
场景 & 推荐写法 & 说明 \\
\midrule
普通变量 & \code{sum}、\code{max_value} & 全小写，单词间用下划线 \\
常量 & \code{MAX_SIZE}、\code{PI} & 全大写，突出“不打算改” \\
结构体类型 & \code{Student}、\code{MotorState} & 首字母大写 \\
函数 & \code{compute_sum}、\code{read_sensor} & 小写 + 动词开头 \\
\bottomrule
\end{tabular}
\end{center}

\begin{pitfall}[这些名字不要用]
\code{int}、\code{float}、\code{for}、\code{while}、\code{return} 是关键字；\code{printf}、\code{scanf}、\code{main} 是标准库里的名字。用它们当变量名，轻则编译报错，重则行为诡异。
\end{pitfall}

\section{基本数据类型}

C 语言的基本类型不多，但每一个都有自己的“大小”和“取值范围”。

\subsection{整数类型}

\begin{center}
\begin{tabular}{@{}llll@{}}
\toprule
类型 & 典型大小 & 典型范围 & 用途 \\
\midrule
\code{char} & 1 字节 & $-128 \sim 127$ & 字符、很小的整数 \\
\code{unsigned char} & 1 字节 & $0 \sim 255$ & 字节数据、位操作 \\
\code{short} & 2 字节 & $-32768 \sim 32767$ & 节省内存的小整数 \\
\code{int} & 4 字节 & 约 $\pm 21$ 亿 & 最常用的整数 \\
\code{unsigned int} & 4 字节 & $0 \sim 4294967295$ & 计数、位掩码 \\
\code{long} & 4 或 8 字节 & 平台相关 & 需要更大范围时 \\
\code{long long} & 8 字节 & 约 $\pm 9.2 \times 10^{18}$ & C99 引入，C89 不一定支持 \\
\bottomrule
\end{tabular}
\end{center}

\begin{note}
“典型大小”只是最常见的情况。C 标准只规定了 \code{char} 一定是 1 字节、\code{short} 不小于 \code{char}、\code{int} 不小于 \code{short} 等相对关系，具体大小由编译器和平台决定。写程序时如果需要精确大小，用 \code{sizeof} 或者 C99 的 \code{int32_t}（C89 下需要自己 typedef）。
\end{note}

\subsection{浮点类型}

\begin{center}
\begin{tabular}{@{}llll@{}}
\toprule
类型 & 典型大小 & 有效数字 & 用途 \\
\midrule
\code{float} & 4 字节 & 约 6\textasciitilde7 位 & 精度要求不高的计算 \\
\code{double} & 8 字节 & 约 15\textasciitilde16 位 & 默认选择，精度更高 \\
\code{long double} & 8/16 字节 & 平台相关 & 特殊高精度场景 \\
\bottomrule
\end{tabular}
\end{center}

\begin{tipbox}[新手先用 \texttt{double}]
除非你明确知道自己在省内存（比如单片机），否则浮点数一律用 \code{double}。\code{float} 的精度比你想的低，累加几次就可能出现明显误差。
\end{tipbox}

\subsection{字符类型}

\code{char} 存的是字符的编码（本质是一个小整数）。例如字符 \code{'A'} 的 ASCII 码是 65：

\begin{codebox}{char.c —— 字符与编码}
#include <stdio.h>

int main(void)
{
    char c = 'A';

    printf("c = %c\n", c);
    printf("c 的编码 = %d\n", c);

    c = c + 1;              /* 字符也可以参与运算 */
    printf("c 加 1 后 = %c\n", c);

    return 0;
}
\end{codebox}

\begin{outbox}{运行结果}
c = A
c 的编码 = 65
c 加 1 后 = B
\end{outbox}

注意：\code{'A'} 是字符常量，用单引号；\code{"A"} 是字符串，用双引号，两者完全不同。字符串的细节第 9 章再讲。

\subsection{字面量怎么写}

字面量就是直接写在代码里的值。

\begin{center}
\begin{tabular}{@{}lll@{}}
\toprule
写法 & 例子 & 说明 \\
\midrule
十进制整数 & \code{42}、\code{-7} & 最常用 \\
八进制整数 & \code{052} & 前导 0，等于十进制 42（不推荐） \\
十六进制整数 & \code{0x2A} & 前缀 \code{0x}，等于十进制 42 \\
无符号后缀 & \code{42U} & 表示 \code{unsigned} \\
长整型后缀 & \code{42L} & 表示 \code{long} \\
浮点数 & \code{3.14}、\code{0.5} & 带小数点 \\
科学计数法 & \code{1.5e3} & 等于 1500.0 \\
字符 & \code{'A'}、\code{'0'}、\code{\n} & 单引号，一个字符 \\
字符串 & \code{"hello"}、\code{""} & 双引号，零个或多个字符 \\
\bottomrule
\end{tabular}
\end{center}

\subsection{转义字符}

有些字符没法直接写出来（比如换行、制表符），或者有特殊含义（比如双引号），就用反斜杠加字母表示：

\begin{center}
\begin{tabular}{@{}lll@{}}
\toprule
转义 & 含义 & ASCII 码 \\
\midrule
\code{\n} & 换行（newline） & 10 \\
\code{\t} & 水平制表符（tab） & 9 \\
\code{\r} & 回车（carriage return） & 13 \\
\code{\b} & 退格 & 8 \\
\code{\\} & 一个反斜杠 & 92 \\
\code{\"} & 双引号 & 34 \\
\code{\'} & 单引号 & 39 \\
\code{\0} & 空字符（字符串结束标志） & 0 \\
\bottomrule
\end{tabular}
\end{center}

\begin{example}
下面这条语句会输出什么？
\begin{codebox}{escape.c —— 转义字符}
#include <stdio.h>

int main(void)
{
    printf("姓名\t年龄\n");
    printf("张三\t18\n");
    printf("他说：\"你好！\"\n");
    return 0;
}
\end{codebox}
\end{example}
\begin{solution}
\code{\t} 输出制表符，对齐两列；\code{\n} 换行；\code{\"} 输出普通的双引号。结果：
\begin{outbox}{运行结果}
姓名    年龄
张三    18
他说："你好！"
\end{outbox}
\end{solution}

\section{输出与输入}

程序要和外界交流：把结果告诉人（输出），把人给的数据读进来（输入）。C 语言的标准库提供了一组函数，最常用的是 \code{printf} 和 \code{scanf}。

\subsection{\texttt{printf}：格式化输出}

\code{printf} 的名字来自 \en{print formatted}（格式化打印）。它的第一个参数是\textbf{格式串}，里面可以嵌入\textbf{占位符}，后面的参数依次填进去。

\begin{center}
\begin{tikzpicture}[x=1cm,y=1cm,>=Stealth]
  \node[font=\ttfamily\small,anchor=west] (f) at (0,0) {printf("a = \%d, b = \%d\textbackslash n", a, b);};
  \draw[ptr] (2.4,-0.25) -- (2.4,-1.1);
  \node[font=\scriptsize,color=ptrline,anchor=north] at (2.4,-1.1) {第 1 个占位符填 \texttt{a}};
  \draw[ptr] (4.6,-0.25) -- (4.6,-1.8);
  \node[font=\scriptsize,color=ptrline,anchor=north] at (4.6,-1.8) {第 2 个占位符填 \texttt{b}};
\end{tikzpicture}
\end{center}

常用的格式符：

\begin{center}
\begin{tabular}{@{}lll@{}}
\toprule
格式符 & 类型 & 例子 \\
\midrule
\texttt{\%d} & 有符号十进制整数 & \texttt{printf("\%d", 42);} \\
\texttt{\%u} & 无符号十进制整数 & \texttt{printf("\%u", 42U);} \\
\texttt{\%x} & 十六进制（小写） & \texttt{printf("\%x", 255);} 输出 \texttt{ff} \\
\texttt{\%o} & 八进制 & \texttt{printf("\%o", 8);} 输出 \texttt{10} \\
\texttt{\%c} & 单个字符 & \texttt{printf("\%c", 'A');} \\
\texttt{\%f} & 浮点数（默认 6 位小数） & \texttt{printf("\%f", 3.14);} \\
\texttt{\%.2f} & 保留 2 位小数 & \texttt{printf("\%.2f", 3.14159);} 输出 \texttt{3.14} \\
\texttt{\%s} & 字符串 & \texttt{printf("\%s", "hi");} \\
\texttt{\%p} & 指针（地址） & \texttt{printf("\%p", (void *)\&a);} \\
\texttt{\%\%} & 输出一个百分号 & \texttt{printf("100\%\%");} \\
\bottomrule
\end{tabular}
\end{center}

\begin{pitfall}[格式符和参数类型必须匹配]
用 \texttt{\%d} 打印 \code{double}，或用 \texttt{\%f} 打印 \code{int}，属于未定义行为。编译器通常会给出警告（\code{-Wall} 打开时），但程序可能仍然“看起来能跑”。请务必让格式符和参数类型一一对应。
\end{pitfall}

\subsection{宽度、精度与对齐}

格式串里还可以指定“占多少格、保留几位小数、左对齐还是右对齐”：

\begin{codebox}{format.c —— 格式化输出}
#include <stdio.h>

int main(void)
{
    printf("[%5d]\n", 42);       /* 宽度 5，右对齐 */
    printf("[%-5d]\n", 42);      /* 宽度 5，左对齐 */
    printf("[%05d]\n", 42);      /* 宽度 5，用 0 填充 */
    printf("[%.3f]\n", 3.14159); /* 保留 3 位小数 */
    printf("[%8.2f]\n", 3.14159);/* 总宽 8，保留 2 位小数 */
    printf("[%c]\n", 'A');
    printf("[%s]\n", "hello");
    return 0;
}
\end{codebox}

\begin{outbox}{运行结果}
[   42]
[42   ]
[00042]
[3.142]
[    3.14]
[A]
[hello]
\end{outbox}

\begin{tipbox}[对齐表格的小技巧]
打印成绩单、实验数据时，用 \texttt{\%8.2f} 这样的固定宽度格式，可以让小数点对齐，输出整齐好看。
\end{tipbox}

\subsection{\texttt{scanf}：读入数据}

\code{scanf} 的名字来自 \en{scan formatted}。它按格式串读取输入，并把结果\textbf{存到变量的地址里}。

\begin{codebox}{input.c —— 读入两个整数}
#include <stdio.h>

int main(void)
{
    int a, b;

    printf("请输入两个整数：");
    scanf("%d%d", &a, &b);

    printf("它们的和是 %d\n", a + b);
    return 0;
}
\end{codebox}

\begin{outbox}{运行示例}
请输入两个整数：3 5
它们的和是 8
\end{outbox}

注意 \code{scanf} 的参数是 \texttt{\&a}、\texttt{\&b}，不是 \code{a}、\code{b}。为什么？

\begin{center}
\begin{tikzpicture}[x=1cm,y=1cm,>=Stealth]
  \node[ptrbox,minimum width=2cm] (a) at (0,0) {\texttt{a}};
  \node[font=\ttfamily\small] (addr) at (2.8,0) {0x1000};
  \node[font=\small,anchor=west] at (4.4,0) {变量的地址};
  \draw[ptr] (a.east) -- (addr.west);
  \node[font=\small,anchor=west] at (-1.2,1.2) {告诉 \texttt{scanf}：“把读到的数写到这个地址”};
\end{tikzpicture}
\end{center}

\texttt{\&} 是“取地址”运算符。\code{scanf} 需要知道把数据写到哪里，所以必须给它地址。关于地址和指针，第 8 章会详细讲，现在先记住“\code{scanf} 的变量前面要加 \texttt{\&}”。

\begin{pitfall}[\texttt{scanf} 的五个坑]
\begin{enumerate}
  \item \textbf{忘记 \texttt{\&}}：\texttt{scanf("\%d", a);} 会导致程序崩溃或读到错误位置。这是最经典的错误。
  \item \textbf{格式串里乱加空格}：\texttt{scanf("\%d, \%d", \&a, \&b);} 要求输入时也必须带逗号。初学者常在格式串里加 \code{\n}，这会让程序等待你多输入一些内容。
  \item \textbf{用 \texttt{\%f} 读 \code{double}}：\code{scanf} 读 \code{double} 必须用 \texttt{\%lf}（注意是字母 l 加 f），而 \code{printf} 里 \texttt{\%f} 和 \texttt{\%lf} 都可以。
  \item \textbf{输入缓冲区残留}：先用 \code{scanf} 读数字，再用 \code{getchar} 读字符时，会读到一个换行符。第 9 章会详细处理。
  \item \textbf{不检查返回值}：输入失败时 \code{scanf} 返回成功读到的项数。健壮的代码应该检查它，例如 \texttt{if (scanf("\%d", \&a) != 1) \{ ... \}}。
\end{enumerate}
\end{pitfall}

\subsection{\texttt{getchar} 与 \texttt{putchar}}

如果只需要读/写单个字符，可以用更简单的 \code{getchar} 和 \code{putchar}：

\begin{codebox}{getchar.c —— 读写单个字符}
#include <stdio.h>

int main(void)
{
    int c;

    printf("请输入一个字符：");
    c = getchar();          /* 从标准输入读一个字符，返回它的编码 */

    printf("你输入的是：");
    putchar(c);
    putchar('\n');

    return 0;
}
\end{codebox}

\begin{outbox}{运行示例}
请输入一个字符：A
你输入的是：A
\end{outbox}

注意这里用 \code{int c} 而不是 \code{char c}：因为 \code{getchar} 在遇到文件结束（EOF）时返回一个特殊值（通常是 $-1$），\code{char} 可能存不下。这个细节第 9 章会展开。

\section{表达式、语句与注释}

\subsection{语句}

语句是程序执行的最小单位。常见的语句有：

\begin{itemize}
  \item 声明语句：\code{int a;}
  \item 表达式语句：\code{a = 10;}、\code{printf("hi");}
  \item 复合语句：用花括号括起来的一组语句 \texttt{\{ ... \}}
  \item 空语句：只有一个分号 \code{;}（有时用于占位）
  \item 控制语句：\code{if}、\code{for}、\code{return} 等（后面几章讲）
\end{itemize}

\subsection{表达式}

表达式是由变量、常量、运算符组成的式子，\textbf{有值}。例如：

\begin{center}
\begin{tabular}{@{}ll@{}}
\toprule
表达式 & 值 \\
\midrule
\code{3 + 5} & 8 \\
\code{a * 2} & \code{a} 的两倍 \\
\code{a = 10} & 赋值表达式的值是 10（赋值本身也有值） \\
\code{printf("hi")} & 打印的字符数（通常用不到） \\
\bottomrule
\end{tabular}
\end{center}

在 C 语言里，\textbf{表达式加上分号就变成语句}。这一点很重要：\code{a + b;} 是合法的语句（虽然没什么用），而 \code{a = b} 是表达式。

\subsection{注释}

注释是写给人看的，编译器会忽略它。C89 只有一种注释：

\begin{codebox}{comment.c —— 注释的写法}
#include <stdio.h>

int main(void)
{
    /* 这是块注释，可以跨行
       也可以写很多行 */
    int a = 10;     /* 行尾注释 */

    /*
     * 这种星号开头的排版只是习惯，不是语法要求
     */
    printf("%d\n", a);
    return 0;
}
\end{codebox}

\begin{c89box}
C89 只支持 \code{/* ... */} 注释，\textbf{不支持} \code{//} 单行注释（那是 C99 引入的）。本书所有代码都遵守这一点。
\end{c89box}

\begin{tipbox}[注释写什么]
好的注释解释“为什么”，而不是“是什么”。\code{a = a + 1;} 不需要注释“让 a 加 1”，但可以注释“记录重试次数”。代码本身应该让人看懂“做什么”，注释用来补充意图、边界条件和坑。
\end{tipbox}

\section{顺序结构综合例题}

顺序结构的程序有一个固定套路：\textbf{输入 → 计算 → 输出}。下面用几个例子把这个套路练熟。

\subsection{例：摄氏度转华氏度}

\begin{example}
输入一个摄氏温度 $C$，输出对应的华氏温度 $F$。公式：$F = C \times \frac{9}{5} + 32$。
\end{example}
\begin{solution}
先看一个想当然的写法：

\begin{codebox}{temperature.c —— 温度转换（错误示范）}
#include <stdio.h>

int main(void)
{
    int c;
    int f;

    printf("请输入摄氏温度：");
    scanf("%d", &c);

    f = c * 9 / 5 + 32;
    printf("华氏温度 = %d\n", f);
    return 0;
}
\end{codebox}

这个程序能跑，但有一个隐藏问题：\code{9 / 5} 是整数除法，结果是 1，不是 1.8。所以 $C=20$ 时算出来是 $20 \times 1 + 32 = 52$，而不是 $68$。

正确写法是把至少一个操作数变成浮点数：
\begin{codebox}{temperature2.c —— 温度转换（正确写法）}
#include <stdio.h>

int main(void)
{
    double c;
    double f;

    printf("请输入摄氏温度：");
    scanf("%lf", &c);

    f = c * 9.0 / 5.0 + 32.0;
    printf("华氏温度 = %.2f\n", f);
    return 0;
}
\end{codebox}

\begin{outbox}{运行示例}
请输入摄氏温度：20
华氏温度 = 68.00
\end{outbox}

要点：\code{9.0 / 5.0} 是浮点除法；\code{scanf} 读 \code{double} 用 \texttt{\%lf}；输出保留两位小数用 \texttt{\%.2f}。
\end{solution}

\subsection{例：交换两个变量}

\begin{example}
输入两个整数 \code{a} 和 \code{b}，交换它们的值后输出。
\end{example}
\begin{solution}
不能直接写 \code{a = b; b = a;}——那样 \code{a} 原来的值会被覆盖。需要借助一个临时变量，就像交换两杯水需要第三个空杯子：

\begin{figure}[htbp]
  \centering
  \begin{tikzpicture}[x=1cm,y=1cm,>=Stealth]
    \node[membox,minimum width=2cm] (a) at (0,0) {a};
    \node[membox,minimum width=2cm] (b) at (2.5,0) {b};
    \node[emptybox,minimum width=2cm] (t) at (5,0) {temp};
    \node[font=\small] at (0,1.0) {3};
    \node[font=\small] at (2.5,1.0) {5};
    \node[font=\scriptsize,color=black!60] at (5,1.0) {?};

    \node[font=\small] at (0,-1.2) {初始};
    \node[font=\small] at (2.5,-1.2) {};
    \node[font=\small] at (5,-1.2) {};

    \draw[ptr] (0,-2.2) to[bend left] node[lbl,below]{① temp = a} (5,-2.2);
    \draw[ptr] (2.5,-2.9) to[bend left] node[lbl,below]{② a = b} (0,-2.9);
    \draw[ptr] (5,-3.6) to[bend left] node[lbl,below]{③ b = temp} (2.5,-3.6);
  \end{tikzpicture}
  \caption{交换两个变量需要借助临时变量}
  \label{fig:swap}
\end{figure}

\begin{codebox}{swap.c —— 交换两个变量}
#include <stdio.h>

int main(void)
{
    int a, b, temp;

    printf("请输入两个整数：");
    scanf("%d%d", &a, &b);

    printf("交换前：a = %d, b = %d\n", a, b);

    temp = a;       /* ① 把 a 存到临时变量 */
    a = b;          /* ② 把 b 写到 a */
    b = temp;       /* ③ 把临时变量写回 b */

    printf("交换后：a = %d, b = %d\n", a, b);
    return 0;
}
\end{codebox}

\begin{outbox}{运行示例}
请输入两个整数：3 5
交换前：a = 3, b = 5
交换后：a = 5, b = 3
\end{outbox}

这个“三步交换”在后面排序算法里会反复出现。
\end{solution}

\subsection{例：三位数反转}

\begin{example}
输入一个三位正整数，输出它的逆序数。例如输入 123，输出 321。
\end{example}
\begin{solution}
利用整数除法和取余：

\begin{itemize}
  \item 个位：\texttt{n \% 10}；
  \item 十位：\texttt{n / 10 \% 10}；
  \item 百位：\code{n / 100}。
\end{itemize}

\begin{codebox}{reverse.c —— 三位数反转}
#include <stdio.h>

int main(void)
{
    int n, ge, shi, bai, rev;

    printf("请输入一个三位数：");
    scanf("%d", &n);

    ge  = n % 10;
    shi = n / 10 % 10;
    bai = n / 100;

    rev = ge * 100 + shi * 10 + bai;
    printf("逆序数 = %d\n", rev);
    return 0;
}
\end{codebox}

\begin{outbox}{运行示例}
请输入一个三位数：123
逆序数 = 321
\end{outbox}

注意 \code{n / 100} 对 123 得到 1（整数除法丢弃小数部分），这正好是百位。
\end{solution}

\subsection{例：求平均值}

\begin{example}
输入三门课的成绩（可能有小数），输出平均分，保留两位小数。
\end{example}
\begin{solution}
思路如下：

\begin{codebox}{average.c —— 求平均分}
#include <stdio.h>

int main(void)
{
    double s1, s2, s3, average;

    printf("请输入三门课成绩：");
    scanf("%lf%lf%lf", &s1, &s2, &s3);

    average = (s1 + s2 + s3) / 3.0;
    printf("平均分 = %.2f\n", average);
    return 0;
}
\end{codebox}

\begin{outbox}{运行示例}
请输入三门课成绩：85.5 90 92.5
平均分 = 89.33
\end{outbox}

注意除以 \code{3.0} 而不是 \code{3}。虽然这里 \code{s1 + s2 + s3} 已经是 \code{double}，除以 \code{3} 也会做浮点除法（因为一个操作数是浮点），但写成 \code{3.0} 更清楚，也更不容易出错。
\end{solution}

\section{常见错误与调试}

初学阶段 90\% 的编译错误都来自下面几类。遇到报错，先对照这张表。

\begin{center}
\begin{tabular}{@{}p{0.34\textwidth}p{0.56\textwidth}@{}}
\toprule
错误现象 & 可能原因 \\
\midrule
\code{expected ';' before ...} & 上一行漏了分号 \\
\texttt{stray '\textbackslash 357' in program} & 用了中文标点（，；：）\\
\code{'printf' undeclared} & 忘了 \code{#include <stdio.h>} \\
\code{undefined reference to 'main'} & 没有 \code{main} 函数，或拼写错误 \\
\code{implicit declaration of function} & 函数没声明/没包含头文件 \\
输出乱码或数字奇怪 & 格式符与类型不匹配，或变量未初始化 \\
程序一闪而过 & 控制台程序运行完自动关闭（正常现象） \\
\bottomrule
\end{tabular}
\end{center}

\begin{tipbox}[读错误信息的顺序]
编译器报错时，\textbf{先看第一条}。后面的错误经常是第一条引起的连锁反应。错误信息里的文件名、行号、列号都要看，尤其是行号。
\end{tipbox}

\begin{deepbox}[拔高：用 \texttt{printf} 调试]
在没有调试器的情况下，最有效的调试手段是“打印中间结果”：
\begin{codebox}{debug.c —— 打印中间变量}
#include <stdio.h>

int main(void)
{
    int a = 10;
    int b = 0;
    int c;

    printf("debug: a=%d, b=%d\n", a, b);
    c = a / b;          /* 这里会出问题 */
    printf("debug: c=%d\n", c);
    return 0;
}
\end{codebox}
运行时会发现程序在除法处崩溃（除零）。通过打印可以看到 \code{b} 是 0。等学了第 4 章的分支，就可以在除法前先判断 \code{b != 0}。
\end{deepbox}

\section{拔高：变量在内存里的细节}

\subsection{声明与定义}

严格来说，“声明”是告诉编译器“有这么个东西”，“定义”是为它分配内存。对于局部变量，\code{int a;} 既是声明也是定义。全局变量和函数的情况稍微复杂一些，等学到多文件编程时再区分。现在你只要知道：\textbf{局部变量在栈上分配，函数返回时自动释放}。

\subsection{为什么局部变量不自动清零}

栈空间会被反复复用。函数 A 用完的栈空间，函数 B 接着用，里面还留着 A 的数据。所以 B 的局部变量如果没有初始化，就会看到“上次留下的垃圾”。这也是为什么“未初始化变量”是危险的。

\subsection{类型转换的初步印象}

当不同类型的值混在一起运算时，C 会自动做类型转换。例如：

\begin{codebox}{convert.c —— 隐式类型转换}
#include <stdio.h>

int main(void)
{
    int a = 7;
    int b = 2;
    double c;

    c = a / b;          /* 先做整数除法，得 3，再转成 3.0 */
    printf("c = %f\n", c);

    c = (double)a / b;  /* 先把 a 转成 double，再做浮点除法 */
    printf("c = %f\n", c);

    return 0;
}
\end{codebox}

\begin{outbox}{运行结果}
c = 3.000000
c = 3.500000
\end{outbox}

\code{(double)a} 叫\textbf{强制类型转换}，第 3 章会系统讲。现在记住：\textbf{整数除法会截断小数部分}，想要精确结果，先把操作数转成浮点。

\section{本章小结}

\begin{center}
\begin{tabular}{@{}p{0.3\textwidth}p{0.6\textwidth}@{}}
\toprule
主题 & 要点 \\
\midrule
程序骨架 & 预处理指令、\code{int main(void)}、函数体、\code{return 0;} \\
语句 & 以分号结尾；花括号组成复合语句 \\
变量 & 名字 + 类型 + 值 + 地址；先声明后使用 \\
C89 声明 & 所有声明必须在块开头，执行语句之前 \\
类型 & \code{char}/\code{int}/\code{float}/\code{double}；大小用 \code{sizeof} 查 \\
输出 & \code{printf} 格式串 + 占位符；格式符与类型必须匹配 \\
输入 & \code{scanf} 要传地址（\texttt{\&a}）；注意返回值与缓冲区 \\
注释 & C89 只用 \code{/* */} \\
顺序结构 & 输入 → 计算 → 输出，自上而下执行 \\
\bottomrule
\end{tabular}
\end{center}

\begin{tipbox}[下一章预告]
顺序结构只能“一条路走到黑”。真实程序需要根据条件做选择（分支），也需要重复做某件事（循环）。在学分支之前，我们先把运算符讲透——尤其是关系运算和逻辑运算，它们是“判断”的基础。
\end{tipbox}

\begin{problemset}
  \item 写一个程序，输出你的姓名、专业和一句座右铭，每行一项。
  \item 输入两个整数，输出它们的和、差、积、商和余数。
  \item 输入一个华氏温度，输出摄氏温度（公式 $C = (F - 32) \times 5 / 9$），保留两位小数。
  \item 输入圆的半径 $r$，输出周长和面积（$\pi$ 取 3.1415926），保留三位小数。
  \item 输入一个四位数，输出它的各位数字之和。例如 2026，输出 10。
  \item 输入两个整数并交换，要求用临时变量实现，并打印交换前后的值。
  \item 下面的程序有几处错误？请找出来并改正。
\begin{codebox}{find-bugs.c —— 找错题}
#include <stdio.h>

int main(void)
{
    int a = 1
    int b = 2;
    sum = a + b;
    printf("%d\n", sum)
    return 0;
}
\end{codebox}
  \item 解释 \texttt{printf("\%d", 3.14);} 为什么是错的，应该怎么写。
  \item 输入一个三位数，输出它的各位数字之和与逆序数。
  \item 思考题：为什么 \texttt{scanf("\%d", \&a);} 需要 \texttt{\&}，而 \texttt{printf("\%d", a);} 不需要？
\end{problemset}
